\begin{exercice}\label{exo1-2}

On consid�re ici une valeur particuli�re pour $f$ :

\[
f_0 (x) = 1,\quad x \in (0,1).
\]

\begin{enumerate}

\item Trouver la solution analytique de \eqref{eq_model} pour $f = f_0$, qu'on notera $u_a$.

\item Soit $u_h$ la solution approch�e calcul�e � l'aide de la m�thode �l�ments finis vue en cours. Calculer explicitement $u_h$ pour une partition uniforme de $(0,1)$ avec $N=1$ point ($h = 1/2$) et $N=2$ points ($h=1/3$).

\item Comparer les deux solutions obtenues $u_h$ � $u_a$, aux points $x=1/3$, $x=1/2$ et $x=2/3$, ainsi qu'en norme $L^1$.\\

\end{enumerate}

%{\bf Equations lin�aires non-homog�nes du 1er ordre.}\\ 

%Trouver la solution des probl�mes de Cauchy suivants :

%\begin{align*}
%&\left\{
%\begin{array}{l}
%\frac{dy}{dx} = 1,\\
%y(0) = 0,
%\end{array}
%\right. && \quad
%%
%\left\{
%\begin{array}{l}
%(1+x^2) \frac{dy}{dx} - x y = 1,\\
%y(0) = 3,
%\end{array}
%\right.\\
%%
%&\left\{
%\begin{array}{l}
%\cos(x) \frac{dy}{dx} + \sin(x) y = x,\\
%y(0) = 1,
%\end{array}
%\right. && \quad
%%
%\left\{
%\begin{array}{l}
%x \frac{dy}{dx} +  y = x,\\
%y(2) = 0.
%\end{array}
%\right.\\
%\end{align*}

%\corrref{TD1_1}

\end{exercice}
